\documentclass[11pt]{article}
\usepackage[margin=1in]{geometry}
\usepackage{amsmath,amssymb}
\newcommand{\finalanswer}[1]{\par\medskip\noindent\fbox{\parbox{0.94\linewidth}{\textbf{Answer:} #1}}}
\title{STAT 412 Homework 5 --- Question 10}
\author{}\date{}
\begin{document}\maketitle
\section*{Problem}
Student heights are approximately normal with $\mu=178.7$ cm and
$\sigma=6.2$ cm. Random samples of size $25$ are drawn. Part (a) asks for the
mean and standard deviation of the sampling distribution of $\bar X$. Part (b)
asks for the expected number, out of 300 sample means recorded to the nearest
tenth, that fall between $177.1$ and $179.7$ cm inclusive. Part (c) asks for
the expected number of sample means falling below $177.0$ cm.

\section*{Work}
For the sampling distribution of a sample mean,
\[
\mu_{\bar X}=\mu=178.7,
\qquad
\sigma_{\bar X}=\frac{\sigma}{\sqrt n}
=\frac{6.2}{\sqrt{25}}=\frac{6.2}{5}=1.24.
\]
Because the means are recorded to the nearest tenth and the interval is
inclusive, use the continuity-corrected endpoints $177.05$ and $179.75$:
\[
z_1=\frac{177.05-178.7}{1.24}\approx -1.33,\qquad
z_2=\frac{179.75-178.7}{1.24}\approx 0.85.
\]
From the z-table,
\[
P(177.1\le \bar X\le 179.7)
\approx \Phi(0.85)-\Phi(-1.33)
=0.8023-0.0918=0.7105.
\]
The expected number is
\[
300(0.7105)=213.15\approx 213.
\]

For part (c), because the sample means are recorded to the nearest tenth,
``below $177.0$'' means recorded values of $176.9$ or lower. Use the class
boundary $176.95$:
\[
z=\frac{176.95-178.7}{1.24}\approx -1.41.
\]
From the z-table,
\[
P(\bar X<176.95)\approx \Phi(-1.41)=0.0793.
\]
Therefore the expected number is
\[
300(0.0793)=23.79\approx 24.
\]

\finalanswer{(a) $\mu_{\bar X}=178.7$ and $\sigma_{\bar X}=1.24$.
(b) Expected number $\approx 213$. (c) Expected number $\approx 24$.}
\end{document}
